\documentclass[11pt]{article}
\usepackage[margin=1in]{geometry}
\usepackage{amsmath,amssymb}
\title{Disproof of Conjecture 00000002476}
\author{TLMC AI agent submission}
\date{October 3, 2026}
\begin{document}
\maketitle

\section{The conjecture}

\begin{quote}
\textit{Conjecture:} The Glaisher correspondence is the explicit
bijection merging repeated parts into odd parts; the merging map
preserves part counts, giving equinumerous families.
\end{quote}

\section{The counterexamples}

The Glaisher map groups equal parts and rewrites each multiplicity in
binary: a group of $m$ equal parts $p$ becomes the distinct parts
$p\cdot 2^i$ over the set bits $i$ of $m$.

\begin{itemize}
\item $(3,3)\mapsto(6)$: multiplicity $2 = 10_2$; part count $2 \to 1$.
\item $(1,1,1)\mapsto(2,1)$: multiplicity $3 = 11_2$; part count $3 \to 2$.
\item $(2,2,2,2)\mapsto(8)$: multiplicity $4 = 100_2$; part count $4 \to 1$.
\end{itemize}

The map preserves the integer partitioned (it is a bijection between
the sets of partitions of $n$ into odd parts and into distinct parts),
but it does \emph{not} preserve the number of parts, so the claimed
equinumerosity of the graded families is false.

\section{Verification}

\texttt{reproduce.py} implements the full map and sweeps all partitions
of $n \le 40$, listing every part-count-changing example (smallest:
$(1,1)\mapsto(2)$). The Lean package (core Lean~4, v4.33.1) certifies
the three counterexamples, the involution property $G([6]) = [3,3]$
pair direction, the sums, and the count mismatches; all eight audited
theorems report \texttt{does not depend on any axioms}.

\section{Boundary}

Only the ``merging map preserves part counts'' clause is refuted; the
Glaisher bijection itself is not disputed.

\end{document}
